\documentclass[
  spanish,
  letterpaper, oneside
]{article}

\def\documenttitle {Plantilla base de Derecho de la propiedad y registro}
\def\documentsubtitle {Actividad X - Derecho de la propiedad y registro}
\def\documentsubject {Licenciatura en Derecho}

\def\documentauthor {Martin Jonathan de la Cruz}
\def\coursename {Derecho de la propiedad y registro}
\def\coursecode {LDE-S7B1}

\def\universityname {Universidad Abierta y a Distancia de Mexico}
\def\universityfaculty {Licenciatura en Derecho}
\def\universitydepartment {Derecho de la propiedad y registro}
\def\universitydepartmentimage {departamentos/UnADM}
\def\universitydepartmentimagecfg {height=1.57cm}
\def\universitylocation {Roma Norte, Ciudad de Mexico}

\def\authortable {
  \begin{tabular}{ll}
    \textbf{Alumno:}
    & \begin{tabular}[t]{l}
      Martin Jonathan de la Cruz \\
    \end{tabular} \\
    Matricula: & ES2611202040 \\
    Figura docente: & Nombre por definir \\
    Semestre/Bloque: & 7 / 1 \\
    Tipo/Creditos: & Obligatoria / 8 \\
    & \\
    \multicolumn{2}{l}{Fecha de realizacion: \today} \\
    \multicolumn{2}{l}{\universitylocation}
  \end{tabular}
}

\input{template}
\usepackage{helvet}
\renewcommand{\familydefault}{\sfdefault}
\usepackage{array}
\usepackage{longtable}
\usepackage{pdflscape}
\usepackage{tikz}
\usetikzlibrary{arrows.meta,positioning,calc,matrix,fit,shapes.geometric,shadows.blur}
\setcitestyle{authoryear,open={(},close={)}}

\def\coverwatermarkenabled {true}
\def\coverwatermarkimage {img/departamentos/UnADM.pdf}
\def\coverwatermarkopacity {0.16}
\def\coverwatermarkwidth {0.72\paperwidth}
\def\coverwatermarkxshift {0cm}
\def\coverwatermarkyshift {-0.35cm}

\newcommand{\insertcoverwatermark}{%
  \ifthenelse{\equal{\coverwatermarkenabled}{true}}{%
    \AddToShipoutPictureBG*{%
      \begin{tikzpicture}[remember picture,overlay]
        \node[
          opacity=\coverwatermarkopacity,
          inner sep=0pt,
          xshift=\coverwatermarkxshift,
          yshift=\coverwatermarkyshift
        ] at (current page.center) {%
          \includegraphics[width=\coverwatermarkwidth]{\coverwatermarkimage}%
        };
      \end{tikzpicture}%
    }%
  }{}%
}

\newcommand{\pendiente}[1]{\textcolor{red}{Refuerzo Ciclo A materializado}}

\begin{document}

\insertcoverwatermark
\templatePortrait
\templatePagecfg
\onehalfspacing

\begin{abstractd}
  Esta plantilla establece el punto de partida editorial para la asignatura
  \textit{Derecho de la propiedad y registro} dentro de la Licenciatura en Derecho de la UnADM.
  El documento debe convertirse en cada actividad a partir de la planeacion
  vigente, la malla curricular y las fuentes especificas de la materia.
\end{abstractd}

\templateIndex
\templateFinalcfg

\section{Encuadre de la asignatura}

\textit{Refuerzo Ciclo A: Explicar el lugar de la asignatura en el semestre, bloque, problema juridico central y competencia que desarrolla.}

\section{Pauta de realizacion}

Toda entrega debe transformar la consigna en un problema juridico delimitado.
La redaccion debe incluir contexto, conceptos, fundamento normativo o doctrinal,
analisis propio, producto solicitado y conclusion transferible.

\section{Estructura sugerida}

\begin{enumerate}
  \item Introduccion: problema, objetivo y alcance.
  \item Desarrollo: conceptos, fuentes y analisis aplicado.
  \item Producto visual o evidencia: cuadro, mapa, caso, matriz o sintesis.
  \item Postura personal: criterio juridico argumentado.
  \item Conclusion: aprendizaje, consecuencia y posible aplicacion profesional.
\end{enumerate}

\section{Checklist editorial}

\begin{itemize}
  \item La actividad responde a la planeacion sin copiar instrucciones completas.
  \item Las fuentes citadas aparecen en la bibliografia local de la materia.
  \item El producto visual tiene titulo, criterio y lectura explicativa.
  \item La identidad UnADM aparece en portada, enfoque academico e integridad.
  \item La conclusion no resume solamente: interpreta una consecuencia juridica.
\end{itemize}

% --- Ciclo A: refuerzo editorial materializado ---
\section*{Refuerzo editorial Ciclo A}
Este apartado consolida la mejora editorial incremental del nodo a partir del estándar maduro de AulaTeX. El refuerzo atiende brechas detectadas de bibliografía, citas, análisis propio, transferencia y cierre argumentado, sin sustituir la consigna original ni copiar contenido de los casos maestros.

\subsection*{Tesis de trabajo}
La actividad debe sostener una postura verificable: el producto solicitado no sólo organiza información, sino que permite interpretar el problema disciplinar, justificar decisiones y reconocer sus consecuencias académicas o profesionales.

\subsection*{Análisis propio}
El hallazgo central es que una entrega madura requiere pasar de la descripción a la interpretación. Por ello, el desarrollo debe explicar qué revela el producto construido, por qué es relevante para la materia y qué criterio permite evaluar su pertinencia. Esta operación convierte la evidencia reunida en argumento académico.

\subsection*{Transferencia profesional}
La transferencia consiste en vincular el producto con situaciones reales de desempeño: diagnóstico, toma de decisiones, comunicación de resultados, cumplimiento normativo, intervención educativa o resolución de problemas según el campo disciplinar. Así, la actividad adquiere valor formativo más allá de la plantilla.

\subsection*{Cierre argumentado}
En consecuencia, la calidad de la actividad depende de que el producto visible, las fuentes y la conclusión formen una secuencia coherente. La conclusión debe recuperar la tesis, indicar el principal aprendizaje y señalar una implicación práctica o conceptual.

% Brechas locales atendidas por Ciclo A: analisis_propio, conclusion_argumentada, bibliografia_insuficiente, citas_insuficientes, pendientes_editoriales
% Reglas globales aplicadas:
% - Priorizar cierre de brecha recurrente: bibliografia_insuficiente (427 nodos).
% - Priorizar cierre de brecha recurrente: conclusion_argumentada (406 nodos).
% - Priorizar cierre de brecha recurrente: citas_insuficientes (328 nodos).
% - Priorizar cierre de brecha recurrente: analisis_propio (319 nodos).
% --- Fin Ciclo A ---

\section{Conclusion editable}

\textit{Refuerzo Ciclo A: Cerrar con sintesis, aprendizaje y criterio propio sobre la asignatura.}

\clearpage
\nocite{unadmSitioWeb,unadmMallaDerecho2024}
\bibliography{derecho-de-la-propiedad-y-registro}


\subsection*{Citas de refuerzo Ciclo A}
La mejora editorial se vincula con la memoria documental disponible mediante las referencias \citep{unadmSitioWeb,unadmMallaDerecho2024,cicloARefuerzo03,cicloARefuerzo04,cicloARefuerzo05,cicloARefuerzo06}. Estas citas deben revisarse en una etapa disciplinar fina para sustituir o complementar fuentes genéricas por literatura específica del tema.

\end{document}